\subsection{Operator morphisms}

\begin{definition}\label{def:operator-morphism}
  Let $(X,L,I)$ and $(Y,M,J)$ be integration systems on sets $X$ and
  $Y$ respectively.
  An \textbf{operator morphism} of integration systems
  \[
    \Lambda:(X,L,I)\to(Y,M,J)
  \]
  is a linear map $\Lambda:L\to M$ satisfying the following two conditions.
  \begin{itemize}
    \item (positivity) for every $f\in L$, if $f\ge0$ then $\Lambda f\ge0$;
    \item (preservation of the integral) $J(\Lambda f)=I(f)$ for every $f\in L$.
  \end{itemize}
\end{definition}

\begin{remark}[relation to the operators of \cite{Yamagami}]\label{rem:yamagami-operator}
  \cite{Yamagami} introduces linear maps $\Lambda$ from a Riesz space $L$ on a
  set $X$ to a Riesz space $M$ on a set $Y$ satisfying
  \begin{itemize}
    \item (positivity) for every $f\in L$, if $f\ge0$ then $\Lambda f\ge0$;
    \item (order continuity) for every sequence of functions
      $\{f_n\}\subset L$, if $f_n(x)\downarrow0$ for all $x\in X$ then
      $(\Lambda f_n)(y)\downarrow0$ for all $y\in Y$.
  \end{itemize}
  Definition~\ref{def:operator-morphism} adds preservation of the integral to
  this and \textbf{removes order continuity}.

  Anything that is an operator in the sense of that book and preserves the
  integral is an operator morphism in the sense of
  Definition~\ref{def:operator-morphism}. The converse does not hold in general.
\end{remark}

\begin{remark}\label{rem:order-continuity-redundant}
  Note that we do not impose the order continuity of
  Remark~\ref{rem:yamagami-operator}.
  \cite{Yamagami} imposes it in order to \emph{construct} an integral
  $J\circ\Lambda$ on $L$ out of $J$ and $\Lambda$. Indeed, if
  $f_n(x)\downarrow0$ then $(\Lambda f_n)(y)\downarrow0$, and since $J$ is an
  integral, $J(\Lambda f_n)\to0$ follows.

  In this paper, by contrast, the integral $I$ is given from the outset as part
  of the object $(X,L,I)$.
  In that situation preservation of the integral, $J\circ\Lambda=I$,
  \textbf{already implies that $J\circ\Lambda$ is an integral}.
  That is, order continuity is unnecessary for the purpose of making
  $J\circ\Lambda$ an integral.

  Order continuity is a genuinely strong condition on the pointwise behaviour of
  $\Lambda f_n$, but it is never used in the discussion below.
  And in fact, for the morphisms of this category, pointwise conditions cannot
  be stated as propositions at all.
  The values of the extension given by
  Theorem~\ref{thm:operator-extension} are elements of the quotient
  $\widetilde{M}$, so a pointwise condition would depend on the choice of
  representative.
  That is, a pointwise condition is not false but meaningless.
  Imposing pointwise conditions at all is inappropriate in this category.

  Note also that the fact that the condition in \cite{Yamagami} is stated in
  terms of pointwise convergence is essential in that book's setting.
  If it were replaced by convergence in the order of $M$ then, since order
  convergence is strictly weaker than pointwise convergence (see the example in
  Appendix~\ref{sec:forget}), $J(\Lambda f_n)\to0$ would no longer follow.
\end{remark}

\begin{lemma}
  \begin{enumerate}
    \item For every integration system $(X,L,I)$, the identity map
      $\operatorname{id}_L:L\to L$ is an operator morphism
      $(X,L,I)\to(X,L,I)$.
    \item For operator morphisms $\Lambda:(X,L,I)\to(Y,M,J)$ and
      $\Xi:(Y,M,J)\to(Z,N,K)$, the composite $\Xi\circ\Lambda:L\to N$ is an
      operator morphism $(X,L,I)\to(Z,N,K)$.
  \end{enumerate}
\end{lemma}
\begin{proof}
  (1) is immediate. As for (2), a composite of linear maps is linear and a
  composite of positive maps is positive.
  Preservation of the integral holds because
  $K(\Xi\Lambda f)=J(\Lambda f)=I(f)$.
\end{proof}

\begin{definition}
  The category whose objects are integration systems and whose morphisms are
  operator morphisms of integration systems is written
  \[
    \mathbf{Int}_{\mathrm{op}} .
  \]
\end{definition}

\subsubsection{Relation to morphisms given by maps of base spaces}

\begin{lemma}\label{lem:base-map-gives-operator}
  Let $F:(X,L,I)\to(Y,M,J)$ be a morphism of integration systems in the sense of
  Section~\ref{chap:preliminaries}.
  Then
  \[
    F^*:L\to M;\quad f\mapsto f\circ\underline{F}
  \]
  is an operator morphism of integration systems.
\end{lemma}
\begin{proof}
  Linearity is immediate from the definition of composition.
  If $f\ge0$ then $(F^*f)(y)=f(\underline{F}(y))\ge0$ for every $y\in Y$, so
  $F^*$ is positive.
  Preservation of the integral is precisely the definition of a morphism of
  integration systems.
  Note that $F^*$ is order continuous as well. Indeed, if $f_n(x)\downarrow0$
  holds for every $x\in X$ then, putting $x:=\underline{F}(y)$ for each
  $y\in Y$, we get $(F^*f_n)(y)=f_n(x)\downarrow0$.
\end{proof}

\begin{proposition}\label{prop:base-to-operator-functor}
  The assignment that leaves objects unchanged and sends a morphism $F$ to
  $F^*$ defines a functor
  \[
    \mathcal{E}:\mathbf{Int}\to\mathbf{Int}_{\mathrm{op}} .
  \]
\end{proposition}
\begin{proof}
  By Lemma~\ref{lem:base-map-gives-operator} morphisms are sent to operator
  morphisms.
  For identity morphisms, $(\operatorname{id}_X)^*=\operatorname{id}_L$.
  As for composition, the map of base spaces underlying $G\circ F$ was
  $\underline{F}\circ\underline{G}$, so for every $f\in L$
  \[
    (G\circ F)^*f
    =f\circ(\underline{F}\circ\underline{G})
    =(f\circ\underline{F})\circ\underline{G}
    =G^*(F^*f) ,
  \]
  and $(G\circ F)^*=G^*\circ F^*$ holds.
\end{proof}

\begin{remark}
  The functor $\mathcal{E}$ need not be faithful.
  Indeed, when $L$ does not separate the points of $X$, distinct maps of base
  spaces $\underline{F}\ne\underline{G}$ can give the same pullback
  $F^*=G^*$.
  In this sense $\mathbf{Int}_{\mathrm{op}}$ does not carry from the outset the
  information about the base space that the integral does not distinguish.
\end{remark}

\begin{remark}
  $\mathcal{E}$ is not full either.
  Among operator morphisms, those coming from maps of base spaces are confined
  to the ones preserving the lattice operations $\wedge,\vee$ (Riesz
  homomorphisms).
  A general operator morphism need not be a lattice homomorphism, and, as
  \cite{Yamagami} points out, this difference is what decides whether
  probability kernels and Markov operators are included.
  The next example exhibits this difference concretely.
\end{remark}

\begin{example}[the case of finite sets: column-stochastic matrices]\label{ex:stochastic-matrix}
  Let $X,Y$ be finite sets and let $L=\mathbb{R}^X$ and $M=\mathbb{R}^Y$.
  Take positive weights $p:X\to(0,\infty)$ and $q:Y\to(0,\infty)$ and put
  \[
    I(f):=\sum_{x\in X}p(x)f(x),
    \qquad
    J(g):=\sum_{y\in Y}q(y)g(y) ;
  \]
  then $(X,L,I)$ and $(Y,M,J)$ are integration systems.

  A linear map $\Lambda:L\to M$ is represented by a matrix $(\Lambda_{yx})$.
  Then
  \begin{itemize}
    \item positivity is equivalent to all entries satisfying
      $\Lambda_{yx}\ge0$;
    \item preservation of the integral, $J(\Lambda f)=I(f)$, says that
      $\sum_{x}\bigl(\sum_{y}q(y)\Lambda_{yx}\bigr)f(x)=\sum_x p(x)f(x)$ holds
      for every $f$, hence is equivalent to
      \[
        \sum_{y\in Y}q(y)\Lambda_{yx}=p(x)
        \qquad(\forall x\in X) .
      \]
  \end{itemize}

  In particular, in the case $p\equiv1$ and $q\equiv1$ (counting, the finite
  version of Example~\ref{ex:counting}), the latter condition says
  \[
    \sum_{y\in Y}\Lambda_{yx}=1
    \qquad(\forall x\in X) ,
  \]
  that is, that \textbf{each column sums to $1$}.
  Hence in this case the operator morphisms of integration systems are precisely
  the \textbf{column-stochastic matrices}.
  Below, ``stochastic matrix'' in this paper means a column-stochastic matrix.

  On the other hand, for a morphism coming from a map of base spaces
  $\underline{F}:Y\to X$ we have $\Lambda_{yx}=1$ when
  $\underline{F}(y)=x$ and $\Lambda_{yx}=0$ otherwise, so each row contains
  exactly one $1$.
  In the counting case, moreover, the column-sum condition forces
  $|\underline{F}^{-1}(x)|=1$, that is, that $\underline{F}$ be a bijection.

  That is, between the same objects,
  \begin{center}
    with maps of base spaces as morphisms only the ``permutation matrices'',\\
    with operators as morphisms all ``column-stochastic matrices''
  \end{center}
  are morphisms. This is the simplest manifestation of the fact that
  $\mathcal{E}$ is not full.
\end{example}

\begin{remark}
  Imposing $\Lambda1_X=1_Y$ (row sums equal to $1$) in the setting of
  Example~\ref{ex:stochastic-matrix} gives what \cite{Yamagami} calls a
  \textbf{Markov operator}, the corresponding family of measures being a
  \textbf{transition probability}.

  Note, however, that in the counting case imposing preservation of the integral
  (column sums equal to $1$) and row sums equal to $1$ simultaneously yields a
  \textbf{doubly stochastic matrix}.
  Hence one cannot say that ``$\mathbf{Int}_{\mathrm{op}}$ contains all
  transition matrices of Markov chains''.
  For general weights $p,q$, the transition matrices that are morphisms are
  confined to those with $\sum_{y}q(y)\Lambda_{yx}=p(x)$, that is, those sending
  $q$ to $p$.
  What is included is exactly those satisfying this condition.
\end{remark}
